\section{A wider strip on the remaining probability panels}

The incoming strip can be extended to the left probability range without
a new Fourier computation. The change is small: retain the negative
conditional amplitude $3/5$, and allow positive amplitude $7/10$.
The following statement is conditional on the same entry and contact
hypotheses as the previously proved strip.

\begin{proposition}\label{left:strip}
Let $n\ge42$ and let $(X,H)$ be an attained fixed-target witness with
nonnegative excess at $C_{\rm E}$ and sharp predecessor
$C_{n-1}\le C_{\rm E}$. Retain the selected v7 coarse entry, its actual
central gap $(-A_-,A_+)$, and the direct one-pass geometry bounds of v10.
For every retained closed probability panel
$[i/100,(i+1)/100]$, $25\le i\le39$, put
\[
 p=\mathbb P(X>0),\quad q=1-p,\quad
 \ell=\frac{\mathbb E X_+}{pq},\quad
 Y=X/\ell+p=B+U,\qquad B=\mathbf1_{\{X>0\}}.
\]
Then the two conditional noises satisfy
\[
 -\frac35<U_0\le a_0,\qquad
 -a_1\le U_1<\frac7{10},\qquad
 a_0,a_1\ge0,\quad a_0+a_1=d<\frac{17}{250}.
\]
In particular $|U_j|<7/10$ on the full left interval $[1/4,2/5]$.
The direct geometric exclusion already gives $p>1/4$ for actual witnesses;
retaining the closed endpoint in the cover only enlarges the tested domain.
\end{proposition}

\begin{proof}
Write $[l,u]$ for a panel and $e,r,b_-,b_+$ for its v10 lower bound on
$\ell$, upper bound on $\rho$, and lower and upper bounds on
$\beta=\mathbb E[X|X|]$. For the negative ray use $b=-b_-$, and for
the positive ray use $b=b_+$. The retained all-threshold heat lemma gives,
for $t>\sqrt3$,
\[
 H_n(H,\pm t)\le h(t),\qquad
 |\partial_t H_n(H,\pm t)|<g(t),
\]
where, with $S=81/80$ and $\eta=8/875$,
\[
 \begin{aligned}
 h(t)&=S^3\left\{\frac25\left(\frac{t^3-3t}{6}
       +\frac{|t^4-6t^2+3|}{72}\right)
       +\frac{t^5}{4100}+\frac{3t}{820}
       +\frac{17(3t^2-2)}{78720}\right\},\\
 g(t)&=\frac25\left(\frac{t^2}{2}+\frac12+\eta\right).
 \end{aligned}
\]
The contact upper estimate is $H_n(H,\pm t)+C_{\rm E}Q(t)$, with
\[
 Q(t)=-t^3+\frac32rt^2+3bt+\eta r.
\]
These statements retain the actual index, nonnegative excess and sharp
predecessor. In particular, they are not consequences of moment bounds
alone. They hold for every raw threshold $H$.

Let $c_0=819/2000<C_{\rm E}$ and $L=g+c_0Q'$. At each proposed
boundary the finite acceptance requires the following eight positive
quantities:
\[
 t^2-3,\quad -h(t)-c_0Q(t),\quad -Q(t),\quad -Q'(t),
 \quad -Q''(t),\quad -L(t),\quad -L'(t),\quad -L''(t).
\]
Here $Q'''=-6$ and $L''=2/5-6c_0=-2057/1000$.
Thus $Q,Q',Q''$ remain negative along the entire outward ray; so do
$L,L'$. The actual contact upper function has derivative strictly less
than $L<0$, because $Q'<0$ and $C_{\rm E}>c_0$. Its boundary value is
negative because $Q<0$ and $h+c_0Q<0$. This excludes the entire ray,
including its boundary. We differentiate the actual heat function using
its derivative estimate, not the absolute-value expression in $h$.

Take
\[
 t_-=e(l+3/5),\qquad t_+=e(1-u+7/10).
\]
The exact rational check passes all eight tests on both rays of all
fifteen closed panels. The weakest negative boundary margin is greater
than $3/10$, and the weakest positive boundary margin is greater than
$27/100$. The input is precisely the source-bound v10 export; no update
is fed back into that export. Shared panel endpoints occur in both
neighbors, so the tests leave no uncovered seam.

It follows that $-t_-<X<t_+$ on the support. Since $\ell\ge e$ and
$l\le p\le u$, division by the actual positive $\ell$ gives
\[
 -3/5<Y<1+7/10.
\]
The actual inner gap gives
\[
 a_0=p-A_-/\ell,\qquad a_1=q-A_+/\ell,\qquad
 a_0+a_1=1-(A_-+A_+)/\ell=d<17/250.
\]
On the negative sign $U_0=Y\le a_0$, while on the positive sign
$U_1=Y-1\ge-a_1$. This proves the asymmetric bounds and hence the
common amplitude claim. Conditional centering is unchanged.
\end{proof}

\paragraph{What the finite calculation does and does not say.}
The bounded exploration considered $0.65,0.70,0.75,0.80,0.85$ and stopped
after the first common success, $0.70$. It made only thirty positive-ray
evaluations. At $0.65$ the boundary test fails on panels 29--33 while the
continuation tests pass. The first failure, on $[.29,.30]$, has boundary
upper estimate strictly between $.010460$ and $.010461$. This is failure
of that sufficient estimate, not a counterexample to the narrower strip.
The acceptance calculation independently reconstructs the expressions at
the selected boundaries; its 31 ray evaluations include this negative
control. No Fourier cells, moment contraction rounds or support-generation
campaign are run.

\section{The resulting count domain}

For the common amplitude $A=7/10$ and displacement ceiling $d_0=17/250$,
\[
 A+2d_0=\frac{209}{250}<1,\qquad
 l_0=\frac{A-1-2d_0}{2}=-\frac{109}{500},\qquad
 h_0=l_0+1=\frac{391}{500}.
\]
The crossing separation is $\zeta=1-d_0=233/250$, and
$-\zeta<l_0<h_0<\zeta$. Every $x=np+H/\ell$ has a shifted anchor
$x=k+u$ with $u\in[l_0,h_0]$. The gap from the amplitude constraints
in the shifted-anchor construction is $(1-A-2d_0)/2=41/500>0$.
Thus the pathwise crossing and curvature lemmas apply on this longer
phase interval.

These anchors are legal. The retained raw-threshold bound is
$|H|\le1.83605680$, and every left-panel scale bound exceeds $49/25$.
Hence $|H/\ell|<1$. For $n\ge42$ and $p\in[1/4,2/5]$,
\[
 np-1-h_0\ge42/4-1-h_0>0,
 \qquad n(1-p)-1+l_0\ge42(3/5)-1+l_0>0.
\]
Consequently $1\le k\le n-1$. This is an exact domain check, not a
bound for the final count comparison. In particular, the variance and
normalization constants proved earlier only for $p\ge.4$ cannot be
reused without a new calculation.

The noise has the same scaling throughout:
\[
 s=\mathbb E U^2=\theta/\ell^2,\qquad
 \lambda=ns=n\theta/\ell^2,\qquad \ell=(pq+s)^{-1/2}.
\]
Together with the existing $3/5$ strips on panels 40--52, the result
supplies a common $7/10$ conditional amplitude on every retained panel
25--52. The signed rays are bounds for one physical law. This statement
does not identify a law with its reflection or replace the original
closed-event reflection and attainment argument.

\subsection{A Bernoulli gap where the external theorem applies}
Schulz's strict finite-$n$ Bernoulli theorem supplies $C_B<C_{\rm E}$
for every $0<p<1$. Its sharper bound applies on $[1/3,2/3]$ and gives
\[
 \rho_B(C_{\rm E}-C_B)
 >\frac{2C_{\rm E}(r-p_*)^2}{\sqrt{pq}},\qquad
 r=\min(p,q),\quad p_*=(4-\sqrt{10})/2.
\]
On the left subinterval $p\in[1/3,2/5]$, both factors
$(p_*-p)^2$ and $(p(1-p))^{-1/2}$ decrease. The minimum is therefore
at $p=2/5$. Using $C_{\rm E}>819/2000$,
$\sqrt{10}<3.162277661$ and $\sqrt{6/25}<49/100$, exact rational
arithmetic gives the simple uniform reserve
\[
 \rho_B(C_{\rm E}-C_B)>\frac{59}{100000}.
\]
Below $p=1/3$ we retain only the global strict Bernoulli bound, with
zero quantitative reserve. Panel 33 must be split at $1/3$ before using
the sharper formula. Neither the sharper formula nor this positive
uniform reserve is asserted below that endpoint.

\subsection{A repair term that changes on the left}
The favorable inequality $\rho\ge\rho_B$ used in the earlier small-noise
argument does not hold on this entire new domain. For an exact example,
take $p=13/50$, $q=37/50$, $U_1=0$ and let $U_0$ take
$\pm\varepsilon$ with equal conditional probabilities, where
$\varepsilon=1/1000$. Put $v=pq$ and $s=q\varepsilon^2$. The centered
raw law has nodes and masses
\[
 (-p-\varepsilon,q/2),\qquad(-p+\varepsilon,q/2),\qquad(q,p),
\]
variance $V=v+s$ and third absolute moment
$M_3=v(1-2v)+3ps$. Exact rational comparison gives
\[
 \frac{M_3^2}{V^3}<\frac{(1-2v)^2}{v},
\]
so the standardized law satisfies $\rho<\rho_B$. Its inward displacement
is only $\varepsilon$, and at $n=42$ its noise is
$\lambda=42s=.00003108<1/500$. The law satisfies the new amplitude and
inward bounds. No contact premise is claimed for this example.
It refutes discarding the adverse moment repair for arbitrary laws in
the enlarged count domain; it does not refute a count comparison.

\paragraph{Remaining obligations.}
The wider strip removes the missing support premise on the left panels.
It does not prove their count inequality. A new argument must control
the own count, incoming counts, normal and moment repairs on the wider
probability and phase ranges. The selected coarse entry, contact and
sharp predecessor remain ancestors of this strip. A finite base is
still needed for the first-failure argument; no result here proves it.
